\section{Interpretación económica de los resultados}

\subsection{La estadística necesita contexto}

Los modelos estadísticos describen patrones observados en datos históricos, pero no conocen por sí mismos las causas económicas de esos patrones. En índices de costos de obras civiles, la interpretación debe integrar el resultado cuantitativo con información del mercado, condiciones contractuales y comportamiento de insumos.

Por ejemplo, una tendencia creciente en el índice de materiales puede estar asociada a inflación general, aumento de demanda de construcción, restricciones de oferta o encarecimiento de importaciones. La regresión identifica la trayectoria, pero el ingeniero debe interpretar si esa trayectoria responde a condiciones persistentes o a un choque transitorio.

\subsection{Acero y materiales metálicos}

El acero es un insumo sensible a precios internacionales, tasa de cambio, costos energéticos y dinámica global de producción. En obras civiles, afecta estructuras metálicas, puentes, refuerzos, elementos prefabricados y componentes de infraestructura. Una serie de índice asociada a acero puede presentar saltos bruscos cuando ocurren restricciones de oferta o incrementos internacionales.

Si el software proyecta una continuidad de crecimiento en una serie relacionada con acero, el informe debe advertir que la validez depende de que no ocurran cambios fuertes en mercados internacionales o políticas comerciales. El modelo captura el patrón histórico local reflejado en el ICOCIV, no todos los determinantes globales del precio.

\subsection{Combustibles, transporte y maquinaria}

Los combustibles inciden en transporte de materiales, operación de maquinaria y costos logísticos. Un aumento del combustible puede transmitirse a múltiples grupos de costo. En términos de serie temporal, este efecto puede generar persistencia: el impacto de un mes puede permanecer durante varios periodos por contratos de transporte, inventarios y ajustes tarifarios.

Cuando una ruta jerárquica del software involucra transporte o maquinaria, la autocorrelación residual puede tener interpretación económica clara. No necesariamente indica un error de datos; puede reflejar rigidez de costos. Aun así, si el modelo no captura esa dependencia, la proyección debe ser interpretada con prudencia.

\subsection{Mano de obra}

La mano de obra puede responder a salarios, regulación laboral, disponibilidad de personal, negociaciones sectoriales y productividad. Sus variaciones pueden ser más graduales que las de insumos importados, aunque también pueden presentar cambios asociados a decisiones normativas o condiciones regionales.

Una serie estable de mano de obra puede ser bien aproximada por un modelo lineal. Sin embargo, una reforma laboral o variación fuerte del salario mínimo puede introducir un cambio no anticipado por los datos históricos. Esta es una limitación propia de modelos univariados basados en tendencia.

\subsection{Inflación y persistencia}

La inflación general puede actuar como fuerza común sobre varios componentes del ICOCIV. Si la economía mantiene inflación positiva, muchos índices de costos muestran crecimiento acumulado. Esto justifica que el software evalúe la coherencia de tendencias no decrecientes en ciertos casos.

No obstante, la inflación no afecta todos los grupos por igual. Materiales importados pueden responder más a la tasa de cambio; insumos locales pueden depender de oferta nacional; transporte puede seguir combustibles. Por ello, la ruta jerárquica seleccionada es crucial para interpretar el resultado.

\subsection{Uso técnico del pronóstico}

El pronóstico del software debe usarse como soporte para análisis, no como sustituto de estudios de mercado. En un informe técnico, una conclusión adecuada podría ser:

\begin{quote}
La serie presenta tendencia creciente y errores predictivos bajos en el backtesting de corto plazo. En consecuencia, la proyección ofrece una referencia cuantitativa razonable para el periodo analizado. Sin embargo, su interpretación debe complementarse con información económica reciente sobre el grupo de costos o insumo seleccionado.
\end{quote}

Esta redacción evita dos extremos: rechazar la utilidad estadística del pronóstico o presentarlo como certeza contractual. El enfoque correcto es integrarlo como evidencia cuantitativa dentro de un análisis profesional más amplio.
